\section{Second-Degree Surfaces in Three-Dimensional Space}

The three-dimensional analogue of a conic section is a \emph{quadric surface}.
Instead of an equation in two coordinates, we now allow a second-degree equation
in three coordinates.

\begin{definitionbox}
A second-degree equation in $\mathbb{R}^3$ has the form
\[
Ax^2+By^2+Cz^2+Dxy+Exz+Fyz+Gx+Hy+Iz+J=0,
\]
where at least one quadratic coefficient is nonzero. Its solution set is called
a quadric surface when it forms a non-degenerate surface.
\end{definitionbox}

Translations move the centre or vertex of the surface. Rotations of the
coordinate system remove mixed terms such as $xy$, $xz$, and $yz$. After these
changes of coordinates, the principal quadric surfaces have simple standard
forms.

\subsection{Ellipsoid}

An ellipsoid centred at the origin has equation
\[
\frac{x^2}{a^2}+\frac{y^2}{b^2}+\frac{z^2}{c^2}=1,
\qquad a,b,c>0.
\]
Its intersections with the coordinate planes are ellipses. When
$a=b=c=r$, the ellipsoid becomes the sphere
\[
x^2+y^2+z^2=r^2.
\]

\begin{center}
\begin{tikzpicture}[scale=0.9]
    \draw[axis] (-3.8,0) -- (4.0,0) node[right] {$x$};
    \draw[axis] (0,-2.8) -- (0,3.0) node[above] {$z$};
    \draw[axis] (0,0) -- (-1.8,-1.4) node[below left] {$y$};
    \draw[subset] (0,0) ellipse (2.8 and 1.7);
    \draw[subset,dashed] (0,0) ellipse (1.0 and 1.7);
    \draw[subset] (0,-1.7) arc[start angle=-90,end angle=90,x radius=1.0,y radius=1.7];
    \draw[subset,dashed] (0,1.7) arc[start angle=90,end angle=270,x radius=1.0,y radius=1.7];
\end{tikzpicture}
\end{center}

\subsection{Elliptic Paraboloid}

An elliptic paraboloid has standard equation
\[
z=\frac{x^2}{a^2}+\frac{y^2}{b^2},
\qquad a,b>0.
\]
Horizontal sections $z=k>0$ are ellipses, while vertical sections through the
$z$-axis are parabolas. The surface has one vertex and opens in one direction.

\begin{center}
\begin{tikzpicture}[scale=0.9]
    \draw[axis] (-3.4,0) -- (3.6,0) node[right] {$x$};
    \draw[axis] (0,-0.5) -- (0,4.6) node[above] {$z$};
    \draw[axis] (0,0) -- (-1.7,-1.3) node[below left] {$y$};
    \draw[subset,domain=-2.2:2.2,samples=100] plot (\x,{0.72*\x*\x});
    \draw[subset] (0,1.1) ellipse (1.25 and 0.42);
    \draw[subset] (0,2.4) ellipse (1.85 and 0.62);
    \draw[subset] (0,3.8) ellipse (2.35 and 0.78);
    \fill[geometricSubsetColor] (0,0) circle (2pt) node[below right] {vertex};
\end{tikzpicture}
\end{center}

\subsection{Hyperbolic Paraboloid}

A hyperbolic paraboloid has standard equation
\[
z=\frac{x^2}{a^2}-\frac{y^2}{b^2}.
\]
Sections parallel to the $xz$-plane and $yz$-plane are parabolas opening in
opposite directions. Horizontal sections are hyperbolas, except at $z=0$, where
the section consists of two intersecting lines. The surface is commonly called
a saddle.

\begin{center}
\begin{tikzpicture}[scale=0.9]
    \draw[axis] (-3.5,0) -- (3.7,0) node[right] {$x$};
    \draw[axis] (0,-2.6) -- (0,3.0) node[above] {$z$};
    \draw[axis] (0,0) -- (-2.0,-1.5) node[below left] {$y$};
    \draw[subset,domain=-2.4:2.4,samples=100] plot (\x,{0.42*\x*\x});
    \draw[subset,domain=-2.2:2.2,samples=100] plot ({0.8*\x},{-0.42*\x*\x});
    \draw[helper] (-2.3,-1.15) -- (2.3,1.15);
    \draw[helper] (-2.3,1.15) -- (2.3,-1.15);
\end{tikzpicture}
\end{center}

\subsection{Hyperboloid of One Sheet}

A hyperboloid of one sheet has standard equation
\[
\frac{x^2}{a^2}+\frac{y^2}{b^2}-\frac{z^2}{c^2}=1.
\]
Horizontal sections are ellipses. Vertical sections through the $z$-axis are
hyperbolas. The surface is connected and has a narrow central waist.

\begin{center}
\begin{tikzpicture}[scale=0.85]
    \draw[axis] (-3.8,0) -- (4.0,0) node[right] {$x$};
    \draw[axis] (0,-3.3) -- (0,3.5) node[above] {$z$};
    \draw[axis] (0,0) -- (-1.8,-1.4) node[below left] {$y$};
    \draw[subset] (0,0) ellipse (1.35 and 0.45);
    \draw[subset] (0,2.2) ellipse (2.15 and 0.72);
    \draw[subset] (0,-2.2) ellipse (2.15 and 0.72);
    \draw[subset] (-1.35,0) .. controls (-1.55,1.0) and (-1.9,1.7) .. (-2.15,2.2);
    \draw[subset] (1.35,0) .. controls (1.55,1.0) and (1.9,1.7) .. (2.15,2.2);
    \draw[subset] (-1.35,0) .. controls (-1.55,-1.0) and (-1.9,-1.7) .. (-2.15,-2.2);
    \draw[subset] (1.35,0) .. controls (1.55,-1.0) and (1.9,-1.7) .. (2.15,-2.2);
\end{tikzpicture}
\end{center}

\subsection{Hyperboloid of Two Sheets}

A hyperboloid of two sheets has standard equation
\[
-\frac{x^2}{a^2}-\frac{y^2}{b^2}+\frac{z^2}{c^2}=1.
\]
The condition forces $|z|\ge c$, so the surface has two disconnected parts.
Horizontal sections are ellipses, while suitable vertical sections are
hyperbolas.

\begin{center}
\begin{tikzpicture}[scale=0.85]
    \draw[axis] (-3.4,0) -- (3.6,0) node[right] {$x$};
    \draw[axis] (0,-4.0) -- (0,4.2) node[above] {$z$};
    \draw[axis] (0,0) -- (-1.6,-1.2) node[below left] {$y$};
    \draw[subset] (0,1.5) ellipse (0.75 and 0.25);
    \draw[subset] (0,3.0) ellipse (1.75 and 0.58);
    \draw[subset] (-0.75,1.5) .. controls (-1.0,2.0) and (-1.45,2.6) .. (-1.75,3.0);
    \draw[subset] (0.75,1.5) .. controls (1.0,2.0) and (1.45,2.6) .. (1.75,3.0);
    \draw[subset] (0,-1.5) ellipse (0.75 and 0.25);
    \draw[subset] (0,-3.0) ellipse (1.75 and 0.58);
    \draw[subset] (-0.75,-1.5) .. controls (-1.0,-2.0) and (-1.45,-2.6) .. (-1.75,-3.0);
    \draw[subset] (0.75,-1.5) .. controls (1.0,-2.0) and (1.45,-2.6) .. (1.75,-3.0);
\end{tikzpicture}
\end{center}

\subsection{Elliptic Cone}

An elliptic cone has equation
\[
\frac{x^2}{a^2}+\frac{y^2}{b^2}-\frac{z^2}{c^2}=0.
\]
It consists of two nappes meeting at the origin. Horizontal sections away from
the origin are ellipses. Vertical sections through the axis are pairs of
intersecting lines.

\begin{center}
\begin{tikzpicture}[scale=0.85]
    \draw[axis] (-3.4,0) -- (3.6,0) node[right] {$x$};
    \draw[axis] (0,-3.5) -- (0,3.7) node[above] {$z$};
    \draw[axis] (0,0) -- (-1.6,-1.2) node[below left] {$y$};
    \draw[subset] (0,0) -- (-2.2,2.8);
    \draw[subset] (0,0) -- (2.2,2.8);
    \draw[subset] (0,0) -- (-2.2,-2.8);
    \draw[subset] (0,0) -- (2.2,-2.8);
    \draw[subset] (0,2.8) ellipse (2.2 and 0.72);
    \draw[subset] (0,-2.8) ellipse (2.2 and 0.72);
\end{tikzpicture}
\end{center}

\subsection{Cylinders as Second-Degree Surfaces}

If one coordinate does not occur in the equation, the corresponding planar
curve is extended parallel to that coordinate axis. For example,
\[
\frac{x^2}{a^2}+\frac{y^2}{b^2}=1
\]
describes an elliptic cylinder parallel to the $z$-axis, while
\[
y=x^2
\]
describes a parabolic cylinder parallel to the $z$-axis.

\begin{center}
\begin{tikzpicture}[scale=0.85]
    \draw[axis] (-3.4,0) -- (3.6,0) node[right] {$x$};
    \draw[axis] (0,-2.7) -- (0,3.0) node[above] {$z$};
    \draw[axis] (0,0) -- (-1.7,-1.3) node[below left] {$y$};
    \draw[subset] (0,1.8) ellipse (1.8 and 0.65);
    \draw[subset] (0,-1.8) ellipse (1.8 and 0.65);
    \draw[subset] (-1.8,1.8) -- (-1.8,-1.8);
    \draw[subset] (1.8,1.8) -- (1.8,-1.8);
    \draw[subset,dashed] (0,1.8) arc[start angle=0,end angle=180,x radius=1.8,y radius=0.65];
    \draw[subset,dashed] (0,-1.8) arc[start angle=180,end angle=360,x radius=1.8,y radius=0.65];
\end{tikzpicture}
\end{center}

\subsection{Degenerate Quadrics}

As in the planar case, a second-degree equation in three variables may be
degenerate. It may describe a point, a line, a plane, two planes, or no real
points. For example,
\[
x^2+y^2+z^2=0
\]
describes only the origin, while
\[
x^2-y^2=0
\]
factors as
\[
(x-y)(x+y)=0
\]
and therefore describes the union of the two planes
\[
x-y=0
\qquad\text{and}\qquad
x+y=0.
\]

\begin{summarybox}
Second-degree equations in three variables produce quadric surfaces. After
translations and rotations, the main non-degenerate types are ellipsoids,
elliptic and hyperbolic paraboloids, hyperboloids of one or two sheets, and
elliptic cones. Equations missing one coordinate give cylinders, while
factorisation may reveal degenerate unions of planes or smaller sets.
\end{summarybox}
